\documentclass[14pt,a4paper]{extarticle}

\usepackage{fontspec}
\usepackage{polyglossia}
\usepackage{geometry}
\usepackage{enumitem}
\usepackage{array}
\usepackage{tabularx}
\usepackage{ragged2e}

\setmainlanguage{russian}
\setmainfont{Times New Roman}

\geometry{
  left=20mm,
  right=15mm,
  top=20mm,
  bottom=20mm
}

\pagestyle{empty}
\setlength{\parindent}{0pt}
\setlength{\parskip}{0pt}
\setlist[enumerate]{leftmargin=*, nosep}

\newcommand{\blankline}{\rule{\linewidth}{0.4pt}}
\newcommand{\shortblank}[1]{\rule{#1}{0.4pt}}
\newcommand{\fieldtext}[1]{\RaggedRight #1}

\begin{document}

\begin{center}
Министерство науки и высшего образования Российской Федерации\\
Санкт-Петербургский политехнический университет Петра Великого\\
Высшая школа прикладной математики и вычислительной физики
\end{center}

\vspace{1.2em}

\begin{flushright}
\begin{minipage}{0.43\textwidth}
УТВЕРЖДАЮ\\
Руководитель ОП\\
Фролов М. Е.\\
<<\shortblank{9mm}>> \shortblank{29mm} 20\shortblank{8mm} г.
\end{minipage}
\end{flushright}

\vspace{1.8em}

\begin{center}
\textbf{ЗАДАНИЕ}\\
на выполнение выпускной квалификационной работы
\end{center}

\vspace{1.1em}

студенту Уртемееву С. А., группа 5030102/20101

\begin{center}
\small фамилия, имя, отчество (при наличии), номер группы
\end{center}

\vspace{0.8em}

\textbf{1. Тема работы:} Расчет загрузки маршрутов общественного транспорта по расписанию.

\vspace{0.7em}

\textbf{2. Срок сдачи студентом законченной работы:} \shortblank{45mm}

\vspace{0.7em}

\textbf{3. Исходные данные по работе:}
\begin{enumerate}
  \item Постановка задачи расчета загрузки маршрутов общественного транспорта по расписанию движения.
  \item Данные о транспортном предложении: остановки, маршруты, рейсы, времена отправления и прибытия транспортных средств.
  \item Данные о пешеходных связях для доступа к остановкам, выхода из системы общественного транспорта и выполнения пересадок.
  \item Данные о пассажирском спросе между транспортными районами с разбиением по временным интервалам.
  \item Параметры модели поиска и отбора соединений: ограничения на пересадки, допустимые времена ожидания, веса компонент обобщенного сопротивления.
  \item Метод назначения пассажирского спроса на сеть общественного транспорта по расписанию с поиском соединений на основе метода ветвей и границ.
  \item Материалы программного проекта и тестовые данные для проверки реализуемой модели.
\end{enumerate}

\vspace{0.7em}

\textbf{4. Содержание работы (перечень подлежащих разработке вопросов):}
\begin{enumerate}
  \item Проанализировать задачу расчета загрузки маршрутов общественного транспорта с учетом расписания движения, пешеходных связей и пассажирского спроса между транспортными районами.
  \item Рассмотреть существующие подходы к назначению пассажирского спроса на сеть общественного транспорта и обосновать применение модели, основанной на расписании движения.
  \item Формализовать математическую модель транспортного предложения, включающую остановки, районы, маршруты, рейсы, времена отправления и прибытия, пешеходные связи, сегменты маршрутов и сегменты соединений.
  \item Формализовать модель пассажирского спроса между транспортными районами с разбиением по временным интервалам.
  \item Разработать модель допустимого пассажирского соединения с учетом времени движения, ожидания, пешеходных перемещений, пересадок, стоимости проезда и обобщенного сопротивления.
  \item Описать алгоритм поиска множества допустимых соединений по расписанию на основе метода ветвей и границ, включая правила продолжения ветвей, критерии доминирования, допуски отбора и ограничения на пересадки.
  \item Разработать процедуру отбора релевантных альтернатив из множества найденных соединений.
  \item Разработать модель распределения пассажирского спроса между выбранными альтернативами с учетом воспринимаемого времени поездки, временной полезности отправления, стоимости проезда и независимости похожих соединений.
  \item Разработать способ расчета загрузок на основе распределенного спроса для транспортных сегментов, элементов рейсов транспортных средств, рейсов, маршрутов и линий.
  \item Разработать способ оценки перегрузки элементов рейсов транспортных средств при заданной провозной способности.
  \item Исследовать возможность расширения алгоритма назначения спроса с учетом провозной способности транспортных средств.
  \item Реализовать расчетный конвейер, включающий загрузку входных данных, предварительную обработку транспортного предложения, поиск соединений, отбор альтернатив, распределение спроса и формирование выходных результатов.
  \item Подготовить тестовые или синтетические данные, необходимые для проверки реализованной расчетной модели.
  \item Выполнить проверочные расчеты, проанализировать корректность распределения спроса и согласованность рассчитанных загрузок.
  \item Сформулировать ограничения разработанной модели и возможные направления ее дальнейшего развития.
\end{enumerate}

\vspace{0.7em}

\textbf{5. Перечень графического материала (с указанием обязательных чертежей):}

Обязательные чертежи не предусмотрены.

Графический материал:
\begin{enumerate}
  \item Схема расчетного конвейера назначения пассажирского спроса на сеть общественного транспорта.
  \item Схема представления транспортного предложения через сегменты маршрутов и сегменты соединений.
  \item Схема поиска пассажирских соединений по расписанию на основе метода ветвей и границ.
  \item Диаграммы результатов проверочного расчета загрузок транспортных сегментов, рейсов, маршрутов и линий.
\end{enumerate}

\vspace{0.7em}

\textbf{6. Перечень используемых информационных технологий, в том числе программное обеспечение, облачные сервисы, базы данных и прочие сквозные цифровые технологии (при наличии):}
\begin{enumerate}
  \item Язык программирования C++20 для реализации основного расчетного конвейера.
  \item Система сборки CMake и менеджер зависимостей vcpkg.
  \item Библиотеки Boost, fmt, spdlog, FTXUI, vincentlaucsb-csv-parser, PEGTL, Eigen и tl-expected.
  \item Локальная библиотека mathfp для типизированных доменных структур, обработки ошибок и вспомогательных математических компонентов.
  \item Язык программирования Python 3.11+ для автономной генерации синтетического пассажирского спроса.
  \item Форматы CSV, JSON и TXT для хранения входных данных, параметров, промежуточных и выходных результатов.
  \item Система контроля версий Git.
\end{enumerate}

\vspace{0.7em}

\textbf{7. Консультанты по работе:} \blankline

\vspace{1.2em}

\textbf{8. Дата выдачи задания} \shortblank{55mm}

\vspace{2.0em}

\begin{tabularx}{\textwidth}{@{}X r@{}}
Руководитель ВКР \shortblank{58mm} & \shortblank{42mm}\\[-0.2em]
\small\hspace{37mm}(подпись) & \small инициалы, фамилия
\end{tabularx}

\vspace{1.7em}

Задание принял к исполнению \shortblank{52mm}

\smallskip
\hspace{57mm}\small (дата)

\vspace{1.4em}

\begin{tabularx}{\textwidth}{@{}X r@{}}
Студент \shortblank{76mm} & Уртемеев С. А.\\[-0.2em]
\small\hspace{23mm}(подпись) & \small инициалы, фамилия
\end{tabularx}

\vfill

\textbf{Примечание:} 1. Данное задание прилагается к ВКР и вместе с ней представляется в ГЭК.

\end{document}
